\documentclass[12pt]{article}
\usepackage{amsmath, amssymb, geometry, setspace}
\geometry{margin=1in}
\setstretch{1.25}
\setlength{\parindent}{0pt}

\title{Solutions -- Quiz 6\\ \vspace{5pt}
\large ECON 320 -- Introduction to Mathematical Economics}
\author{Instructor: Muhebullah Karimzada}
\date{December 01, 2025}

\begin{document}
\maketitle

\section*{Part A — Multiple Choice}

\begin{enumerate}
\item (c)  
\item (b)  
\item (a)  
\item (b)  
\item (c)  
\item (b)  
\item (d)  
\end{enumerate}

\newpage

\section*{Part B — Short Answer Solutions}

\subsection*{Question 1 — Envelope Theorem}

\textbf{(1) Derivative of the value function}

From the Envelope Theorem for a constrained maximization problem:
\[
\boxed{
\frac{dV}{da}
= U_a + \lambda\, g_a
}
\]
evaluated at the optimum.

\textbf{(2) Why indirect effects disappear}

The FOCs ensure that the derivatives with respect to $x$ and $y$ cancel out. Thus changes in $a$ affect $V(a)$ only directly through $U_a$ and indirectly through the constraint term $g_a$.

\textbf{(3) Interpretation of $-\lambda g_a$ (or $+\lambda g_a$ depending on constraint form)}

It measures how changes in the parameter $a$ affect the tightness of the constraint, multiplied by the shadow value of relaxing the constraint. It is the effect on the optimized value caused by changes in the feasible set.

---

\subsection*{Question 2 — Duality and Shephard’s Lemma}

\textbf{(1) Conditional factor demands}

By Shephard’s Lemma:
\[
\boxed{
\frac{\partial C}{\partial w} = L^*, \qquad
\frac{\partial C}{\partial r} = K^*.
}
\]

\textbf{(2) Duality explanation}

Cost minimization and output maximization share the same tangency condition ($MRTS = w/r$). Therefore the cost function and production function describe the same technology from two perspectives: one minimizes cost for a given output, the other maximizes output for given inputs.

\textbf{(3) Marginal cost}

\[
\boxed{
MC(Q) = \frac{\partial C}{\partial Q}
}
\]
because the Lagrange multiplier in the cost-minimization problem equals marginal cost.

---

\subsection*{Question 3 — Capital Dynamics}

The differential equation is:
\[
\frac{dK}{dt}=I(t)-\delta K.
\]

\textbf{Solution:}
\[
\boxed{
K(t)
=K(0)e^{-\delta t}+\int_0^t e^{-\delta (t-s)}I(s)\,ds.
}
\]

\textbf{Steady state (constant $I$):}
\[
\boxed{
K^*=\frac{I}{\delta}.
}
\]

\textbf{Interpretation:}
- $K(0)e^{-\delta t}$ = initial capital depreciated over time.  
- The integral term = accumulated past investment adjusted for depreciation.

---

\subsection*{Question 4 — Present Value}

Income:
\[
Y(t) = Y_0 e^{gt}.
\]

Present value:
\[
PV = \int_0^\infty Y_0 e^{(g-r)t} dt
= \frac{Y_0}{r-g},
\quad r>g.
\]

\textbf{Condition for finiteness:}  
\[
\boxed{r>g}
\]
because income must not grow faster than it is discounted.

\end{document}
